% !TeX program = xelatex
% !TeX encoding = UTF-8
% 清华大学专业学位硕士研究生选题报告及论文工作计划（202401版）
% 独立文件：在 Overleaf 中将本文件设为 Main document，选择 XeLaTeX。
% 依据用户提供的 Word 表格转制；默认保留原件四页 B5 结构。
\documentclass[UTF8,fontset=fandol,zihao=5]{ctexart}
\usepackage{iftex}
\ifXeTeX\else
  \PackageError{imdt}{Please select XeLaTeX as the compiler}{This template requires XeLaTeX.}
\fi
\usepackage[paperwidth=182mm,paperheight=257mm,left=15mm,right=15mm,
  top=28mm,bottom=20mm]{geometry}
\usepackage{fontspec}
\setmainfont{texgyretermes-regular.otf}[
  BoldFont=texgyretermes-bold.otf,
  ItalicFont=texgyretermes-italic.otf,
  BoldItalicFont=texgyretermes-bolditalic.otf]
% 中文默认 Fandol 宋体/黑体；不依赖本机字体或外部文件。
% 若院系另有字库要求，可合法上传字体后替换：
% \setCJKmainfont{SimSun} \setCJKsansfont{SimHei}
% \setmainfont{Times New Roman}
\usepackage{amssymb}
\usepackage{array,tabularx}
\usepackage{enumitem}
\usepackage{xcolor}
\usepackage[hidelinks,unicode]{hyperref}

% ==================== 个人信息填写区 ====================
% 两个 .tex 独立编译，个人信息请分别填写。
\newcommand{\StudentName}{}
\newcommand{\StudentID}{}
\newcommand{\SchoolName}{深圳国际研究生院}
\newcommand{\MajorName}{电子信息（互动媒体设计与技术）}
\newcommand{\Supervisor}{}
\newcommand{\SupervisorRank}{}
\newcommand{\EnrollmentDate}{}
\newcommand{\FormDate}{\hspace{1em}年\hspace{1em}月\hspace{1em}日}
\newcommand{\ThesisTitle}{}
\newcommand{\TopicScope}{}
\newcommand{\ClassificationNumber}{}
\newcommand{\ProposalMeetingDate}{\hspace{1em}年\hspace{1em}月\hspace{1em}日}
\newcommand{\CrossDisciplineName}{}
\newcommand{\DomesticReferenceCount}{}
\newcommand{\ForeignReferenceCount}{}

% 需要勾选的项目，把对应的 {} 改为 {\tick}；未选项目留空。
% 每组“是/否”只能勾选一项；论文类型只可勾选一项。
\newcommand{\PublicMark}{}
\newcommand{\SecretMark}{}
\newcommand{\ConfidentialMark}{}
\newcommand{\CrossYes}{}
\newcommand{\CrossNo}{}
\newcommand{\ExpertYes}{}
\newcommand{\ExpertNo}{}
\newcommand{\ExperimentYes}{}
\newcommand{\ExperimentNo}{}
\newcommand{\TrainingYes}{}
\newcommand{\TrainingNo}{}
\newcommand{\SafetyReportYes}{}
\newcommand{\SafetyReportNo}{}
\newcommand{\TypeResearch}{}
\newcommand{\TypeSurvey}{}
\newcommand{\TypeCase}{}
\newcommand{\TypeProduct}{}
\newcommand{\TypeScheme}{}
\newcommand{\TypeOther}{}
\newcommand{\OtherTypeDescription}{}

% 可在下列宏中直接填写文字；段落之间可以使用 \par。
\newcommand{\KeyReferences}{}
\newcommand{\AdvisorComment}{}
\newcommand{\ResearchConditions}{}
\newcommand{\ProjectCategory}{}
\newcommand{\MainWorkPlan}{}
\newcommand{\LiteratureCompletionDate}{}
\newcommand{\ResearchCompletionDate}{}
\newcommand{\WritingCompletionDate}{}
% 考核意见、成绩和手写签名栏按原模板保留空白，交相应人员填写。

\newif\ifShowFormHints
\ShowFormHintstrue
% 正式打印时可改为 \ShowFormHintsfalse，隐藏原件中的红色填写提示。
\newcommand{\formhint}[1]{\ifShowFormHints{\color{red}#1}\fi}

% ==================== 表格格式（一般无需修改） ====================
\pagestyle{empty}
\setlength{\parindent}{0pt}
\setlength{\parskip}{0pt}
\setlength{\tabcolsep}{4pt}
\setlength{\arrayrulewidth}{0.4pt}
\linespread{1}
\xeCJKDeclareCharClass{CJK}{"2460 -> "2465}
\AtBeginDocument{\fontsize{10.5bp}{14bp}\selectfont}
\newcommand{\tick}{\ensuremath{\checkmark}}
\newcommand{\option}[1]{（\makebox[1em][c]{#1}）}
\newcommand{\field}[2][28mm]{\underline{\makebox[#1][c]{#2}}}
\newcommand{\blankdate}{\hspace{1em}年\hspace{1em}月\hspace{1em}日}
\newlength{\FormInnerWidth}
\setlength{\FormInnerWidth}{\dimexpr\textwidth-2\tabcolsep-2\arrayrulewidth\relax}
% 单元格设置“最小高度”，填写内容较多时会自然增高，不裁切文字。
% 每页表格不自动跨页；超出预留空间时请缩减文字或另附续页。
\newcommand{\FormCell}[2]{%
  \begingroup
  \setbox0=\vbox{%
    \hsize=\linewidth
    \fontsize{10.5bp}{14bp}\selectfont
    \setlength{\parindent}{0pt}\setlength{\parskip}{0pt}%
    \vspace{1.5mm}#2\par\vspace{1.5mm}}%
  \dimen0=#1\relax
  \ifdim\dimexpr\ht0+\dp0\relax<\dimen0
    \setbox0=\vbox to \dimen0{\unvbox0\vfil}%
  \fi
  \leavevmode\box0
  \endgroup}
\newenvironment{FormTable}
  {\par\noindent\begin{tabular}{|p{\FormInnerWidth}|}\hline}
  {\end{tabular}\par}
\setlist[enumerate]{label=\arabic*.,leftmargin=2em,itemsep=0pt,
  topsep=0pt,parsep=0pt,partopsep=0pt}
\hypersetup{pdftitle={研究生选题报告及论文工作计划},pdfauthor={}}

\begin{document}

% ==================== 第 1 页：封面 ====================
\begingroup
\centering
\vspace*{11mm}
{\songti\fontsize{22bp}{30bp}\selectfont 清\hspace{1.1em}华\hspace{1.1em}大\hspace{1.1em}学\par}
\vspace{12mm}
{\songti\fontsize{18bp}{26bp}\selectfont 攻\hspace{0.4em}读\hspace{0.4em}硕\hspace{0.4em}士\hspace{0.4em}学\hspace{0.4em}位\par}
\vspace{6mm}
{\songti\fontsize{18bp}{26bp}\selectfont 研究生选题报告及论文工作计划\par}
{\fontsize{12bp}{20bp}\selectfont （202401版）\par}
\vspace{23mm}
\fontsize{12bp}{20bp}\selectfont
\begin{tabular}{@{}l@{}}
研究生姓名\ \field[36mm]{\StudentName}\quad 学号\ \field[39mm]{\StudentID}\\[9mm]
院\hspace{3em}系\ \field[99mm]{\SchoolName}\\[9mm]
专业学位类别\ \field[91mm]{\MajorName}\\[9mm]
指导教师\ \field[40mm]{\Supervisor}\ 专业技术职务\ \field[28mm]{\SupervisorRank}\\[9mm]
入学日期\ \field[99mm]{\EnrollmentDate}
\end{tabular}
\vfill
{\fontsize{12bp}{20bp}\selectfont \FormDate\par}
\vspace*{8mm}
\endgroup
\clearpage

% ==================== 第 2 页：选题与文献 ====================
\begin{FormTable}
  \FormCell{51mm}{%
    论文题目：\field[117mm]{\ThesisTitle}\par
    或选题范围：\field[113mm]{\TopicScope}\par
    论文密级请在（\quad）中打“\tick”：公开\option{\PublicMark}\quad
    秘密\option{\SecretMark}\quad 机密\option{\ConfidentialMark}\textsuperscript{1}\par
    定密申请表编号：\field[43mm]{\ClassificationNumber}
    {\fontsize{9bp}{12bp}\selectfont （对于密级为秘密或机密的学位论文）}\par
    选题报告会日期：\ProposalMeetingDate\par
    论文选题是否涉及跨专业学位类别：是\option{\CrossYes}\textsuperscript{2}
    \field[31mm]{\CrossDisciplineName}\quad 否\option{\CrossNo}\par
    {\fontsize{9.5bp}{13bp}\selectfont
    若涉及，选题报告会是否已聘请所涉其他专业学位类别的专家参加：
    是\option{\ExpertYes}\quad 否\option{\ExpertNo}\par}
  }\\\hline
  \FormCell{56mm}{%
    选题报告要求：
    \begin{enumerate}
      \item 选题报告的字数一般不少于五千字；
      \item 选题报告的内容应包括：选题意义、该领域国内外研究动态（文献综述）、本课题研究的目的、预期成果、研究方案、研究方法及其论证、关键难点拟采取的解决措施、论文工作总体日程安排，预计答辩时间等；
      \item 要求查阅一定数量的中、外文献资料。文献综述部分不是将文献内容进行简单的堆砌，而应通过阅读、消化、提炼，对已有的研究成果和动态进行全面的综述；
      \item 选题报告的考核由两部分组成：书面报告和口头报告，分别评分后给出一个总成绩；
      \item 选题报告和研究工作应具有正确的政治立场，符合学术规范；
      \item 填好“研究生选题报告及论文工作计划”表格后，连同书面报告一起交院系研究生管理部门备案。
    \end{enumerate}
  }\\\hline
  \FormCell{78mm}{%
    阅读国内外文献情况：\par
    \hspace{1em}国内文献约\field[22mm]{\DomesticReferenceCount}篇，
    国外文献约\field[26mm]{\ForeignReferenceCount}篇。\par
    主要文献：\par
    \KeyReferences
  }\\\hline
\end{FormTable}
\vspace{2mm}
{\fontsize{8bp}{10.5bp}\selectfont
\textsuperscript{1} 涉密学位论文的定密申请应当在论文开题之前进行。\par
\textsuperscript{2} 如选“是”，应按国家现行《研究生教育学科专业目录》规范填写所涉其他专业学位类别的名称。\par}
\clearpage

% ==================== 第 3 页：评语、考核与签名 ====================
\begin{FormTable}
  \FormCell{76mm}{%
    导师评语：（就论文选题意义、研究生对国内外研究现状的了解情况、研究方法、研究手段、预期成果、学术规范等予以评价）\par
    \vspace{6mm}
    \formhint{须有导师评语、导师手写签名及日期}\par
    \AdvisorComment
    \vfill
    \hfill 导师签名：\field[25mm]{}\quad 日期：\blankdate\par
  }\\\hline
  \FormCell{76mm}{%
    选题报告会考核成员组意见：（选题和研究工作有无政治立场问题、选题意义、论文工作是否具备条件、是否同意选定该课题；若选题涉及跨专业学位类别，考核小组是否已包含所涉其他专业学位类别的专家，选题是否仍符合所属专业学位类别的要求等）\par
    \vspace{8mm}
    \formhint{不填}\par
    \vfill
    \hspace{2em}考核成员组签名：\par
    \vspace{5mm}
  }\\\hline
  \FormCell{44mm}{%
    选题报告考核成绩：（以等级制方式记载。以下评分方法供参考，鼓励院系结合专业学位类别特点制定评分规则。）\par
    书面报告成绩：\field[32mm]{\formhint{导师填写}}（占60\%，由导师评定）\par
    口头报告成绩：\field[32mm]{}（占40\%，由考核成员组评定）\par
    总成绩：\field[28mm]{}\quad 研究所（室）主任签名：\field[25mm]{}\par
    \hfill 日期：\blankdate\par
  }\\\hline
\end{FormTable}
\clearpage

% ==================== 第 4 页：论文工作计划 ====================
{\centering\heiti\fontsize{14bp}{20bp}\selectfont
论\quad 文\quad 工\quad 作\quad 计\quad 划\par}
\vspace{2mm}
\begin{FormTable}
  \FormCell{23mm}{%
    实验设备或其他研究条件落实情况：\par
    \ResearchConditions
  }\\\hline
  \FormCell{23mm}{%
    实验安全落实情况请在（\quad）中打“\tick”：\par
    （1）论文工作中是否涉及实验操作：\hfill 是\option{\ExperimentYes}\quad 否\option{\ExperimentNo}\hspace{12mm}\par
    （2）是否参加安全教育培训并通过考试：\hfill 是\option{\TrainingYes}\quad 否\option{\TrainingNo}\hspace{12mm}\par
    （3）是否已提交实验安全分析报告\textsuperscript{3}：\hfill 是\option{\SafetyReportYes}\quad 否\option{\SafetyReportNo}\hspace{12mm}\par
  }\\\hline
  \FormCell{28mm}{%
    课题所属科研项目类别（国家项目、省部委等政府部门项目、企业委托研究项目、学校项目、自选题目、或其他）：\par
    \ProjectCategory
  }\\\hline
  \FormCell{6mm}{论文类型：在所属类型前划“\tick”（只可选择一项填写）}\\\hline
  \FormCell{17mm}{%
    \begin{tabularx}{\linewidth}{@{}p{5mm}X|p{5mm}X@{}}
      \TypeResearch & ① 专题研究类论文 & \TypeSurvey & ② 调研报告\\\hline
      \TypeCase & ③ 案例分析报告 & \TypeProduct & ④ 产品设计（作品创作）\\\hline
      \TypeScheme & ⑤ 方案设计 & \TypeOther & ⑥ 其它（详细注明）\OtherTypeDescription
    \end{tabularx}
  }\\\hline
  \FormCell{79mm}{%
    论文工作主要内容及日程安排：\par
    \MainWorkPlan
    \vfill
    文献阅读、科研调查计划完成日期：\field[81mm]{\LiteratureCompletionDate}\par
    论文实际工作预计完成日期：\field[91mm]{\ResearchCompletionDate}\par
    论文撰写预计完成日期：\field[98mm]{\WritingCompletionDate}\par
  }\\\hline
\end{FormTable}
\vspace{1mm}
本材料留院系研究生管理部门供中期检查、论文抽查时参考。\par
\vspace{1mm}
{\fontsize{8bp}{10.5bp}\selectfont
\textsuperscript{3} 实验安全分析报告使用清华大学实验室管理处网站模板。\par}

\end{document}
